\section{Non-Normative Reference Profile Score}

This section is non-normative. The Jankurai standard is stack-neutral: a Go, JVM, .NET, Python, Rails, Elixir, or TypeScript-heavy repository can conform if it emits equivalent owner routing, proof receipts, generated-zone evidence, and merge witnesses. The Reference Profile Score below is an aid for durable product systems maintained by agents and humans together. It deliberately weights correctness, security, contracts, and repairability above first-draft velocity because AI already lowers the cost of plausible drafts. The weights in this edition are policy assumptions, not universal empirical truth.

\textbf{The scoring model is not required to implement Jankurai Core and MUST NOT reject an otherwise conformant alternate profile.}

\begin{table*}[t]
\caption{Non-normative Jankurai reference-profile scoring rubric.}
\label{tab:stack-rubric}
\centering
\small
\begin{tabularx}{\textwidth}{L{0.21\textwidth} R{0.07\textwidth} Y Y}
\toprule
\JKTableHeader
\textbf{Criterion} & \textbf{Weight} & \textbf{Full-credit signal} & \textbf{Evidence and argument} \\
\midrule
Agent-verifiable correctness loop & 18 & Fast type/compile checks, deterministic test routing, generated contracts, property tests, reproducible local and CI lanes. & AI output must be treated as suspect until a proof loop rejects or accepts it; setup and context benchmarks show environment and retrieval reliability are central bottlenecks \cite{setupBench2025,contextBench2026}. \\
Security, memory safety, supply chain & 16 & Memory-safe core, secret scanning, SCA, SBOM/provenance, lockfiles, unsafe ledger, dependency rationale. & CISA/NSA memory-safety guidance, Veracode GenAI findings, and GitGuardian secrets evidence all push security into structural gates \cite{cisaMemorySafeLanguages2025,veracodeGenaiCodeSecurity2025,gitguardianSecretsSprawl2026}. \\
Runtime performance, memory, cloud cost & 13 & Efficient hot path, predictable latency, bounded memory, cheap concurrency, low cold-start cost. & Agent-era systems are compute-heavy and parallel; runtime cost becomes architectural, not cosmetic. \\
Concurrency and resilience & 12 & Structured async, cancellation, backpressure, idempotency, retries, dead-letter handling, replayable workflows. & Go survey evidence shows broad AI use but persistent quality concerns; concurrency needs guardrails rather than cleverness \cite{goDeveloperSurvey2025}. \\
Boundary and contract integrity & 11 & OpenAPI/Protobuf/JSON Schema sources, generated clients, schema drift checks, migration gates. & Cross-language systems work only when contracts are generated and tested; OpenAPI Generator and Buf provide mature paths \cite{openapiGeneratorUsage2026,bufDocs2026}. \\
Simplicity and review surface & 10 & Small files, narrow modules, explicit ownership, boring dependency direction, low duplication. & Agents over-edit when boundaries are implicit; SWE-agent and SWE-bench show interfaces and environment shape affect agent results \cite{sweAgentAci2024,sweBench2024}. \\
Ecosystem, hiring, AI corpus strength & 8 & Large corpus, maintained packages, common deployment patterns, strong model familiarity. & GitHub Octoverse language patterns and TypeScript's rise matter because public examples are part of the agent substrate \cite{githubOctoverse2025}. \\
Observability, auditability, provenance & 6 & OpenTelemetry traces/metrics/logs, request IDs, release provenance, repair receipts. & OpenTelemetry gives a vendor-neutral vocabulary for post-merge repair \cite{openTelemetryDocs2025,otelSemanticConventions2026}. \\
Data integrity and durable workflow fit & 4 & PostgreSQL constraints, migrations, indexes, RLS where useful, transactional workflow safety. & PostgreSQL constraints and row-security policies move durable truth below fallible application code \cite{postgresConstraints2026,postgresRowSecurity2026}. \\
Product velocity and interop & 2 & Fast UI loop, generated clients, easy local dev, standard browser ecosystem. & TypeScript 7 and Vite 8 improve feedback speed, but velocity is deliberately weighted below proof \cite{typescript70Beta2026,vite80Release2026}. \\
\bottomrule
\end{tabularx}
\end{table*}

Points are awarded only for structural properties. ``We review carefully'' is not a security control. ``The senior engineer knows where SQL belongs'' is not a boundary. ``The README explains it'' is not a generated contract. A reference profile should make the wrong path hard to express and the right path cheap to prove.

Hard filters apply before points. A profile with no deterministic local proof lane, no security lane for high-risk changes, no generated contract discipline, or no credible data boundary cannot be recommended even if it is pleasant to write. These caps are policy bands and should be recalibrated against incident data, repair time, false positives, and review burden.

\begin{table}[t]
\caption{Representative hard caps before weighted scoring.}
\label{tab:hard-caps}
\centering
\scriptsize
\begin{tabularx}{\columnwidth}{Y R{0.18\columnwidth}}
\toprule
\JKTableHeader
\textbf{Missing or unsafe control} & \textbf{Max score} \\
\midrule
No deterministic local proof lane & 65 \\
No security lane on high-risk repo & 60 \\
No generated contract discipline for public boundary & 80 \\
Direct DB access from wrong layer & 66 \\
No secret or dependency scanning in CI & 78 \\
Prompt injection or overbroad agent agency in trusted policy & 65--78 \\
No rendered UX proof for user-facing web surface & 84 \\
False-green tests or destructive migration without safety evidence & 70--76 \\
\bottomrule
\end{tabularx}
\end{table}
